% Notebook 03, Parts 7 and 8. English report fragment.
% Requires booktabs and graphicx. Include the companion appendix fragment once.
\subsection{Static plans under uncertain future conditions}
\label{sec:nb03-static-comparison}

The baseline, East stretch (Static~1) and full route (Static~2) were evaluated under 2,000 identical Latin-hypercube scenarios, giving 6,000 plan evaluations. The saved transport surrogate provides Year-1 and Year-40 operating snapshots; these comparisons do not simulate adaptive construction decisions. Cycling shares count trips, including walking in the denominator, rather than distance-weighted passenger-kilometres.

Year-40 cycling trip shares have medians of 22.01\%, 22.11\% and 22.21\%, respectively, with p10--p90 ranges of 17.12--28.42\%, 17.21--28.52\% and 17.30--28.64\%. The distributions overlap substantially: uncertainty in future conditions is much larger than the infrastructure effect. Road-delay and operating-emission distributions also overlap (Appendix~\ref{app:nb03-plan-evidence}). Same-future comparisons nevertheless show consistent cycling increases and road-delay reductions. Median cycling increases are 0.096 and 0.205 percentage points for the East stretch and full route; median annual delay reductions are 6,530 and 10,357 person-hours. Operating CO$_2$ changes are less consistent: median changes are +12.2 and $-8.4$ tonnes/year, with p10--p90 ranges of $-4.5$ to +39.7 and $-33.1$ to +13.1 tonnes/year. These are annual car-and-rail operating emissions, excluding construction.

For cycling trip share, feature scores identify cycling preference as the dominant predictive driver in every plan (approximately 95.5--95.6\%), followed by e-bike share (2.8--2.9\%). Rankings barely change after construction. These scores describe predictive importance, not causal effects or shares of explained variance. Monitoring cycling preferences and e-bike adoption is therefore relevant across all plans. Neither cycling gains nor operating emissions alone establish a preferred plan; lifecycle costs, construction emissions and other benefits require later appraisal.

\subsection{Performance relative to a model-based benchmark}
\label{sec:nb03-vulnerability}

A Year-40 cycling trip share of 22\% is used as a model-based performance benchmark, obtained by rounding the baseline median of 22.0085\%. It is an internal reference for a typical modelled baseline future, not an official policy target, external forecast or independently justified minimum ambition. All plans use the same threshold and count walking trips in the denominator. A scenario is classified as below-benchmark when its cycling trip share is strictly less than 22\%; equality meets the benchmark. About half the baseline sample is expected to fall below a median-based threshold, which is not independent evidence of inadequate performance.

The below-benchmark fractions are 49.90\% for the baseline (998 cases), 48.85\% for the East stretch (977), and 47.65\% for the full route (953), out of 2,000 scenarios each. Both classes exceed the minimum of 15 cases. Infrastructure therefore modestly reduces the incidence of below-benchmark outcomes in the shared sample; these fractions are not real-world probabilities.

PRIM searches the five transport uncertainties, excluding appraisal factors that cannot alter cycling trip share. Every plan uses the same 75\% minimum density and must restrict at least one input. The selected baseline box combines the Year-40 input ranges listed in Appendix~\ref{app:nb03-prim-evidence}; the two-dimensional plot displays cycling preference and e-bike share but omits its additional PT-preference restriction. These joint conditions describe relatively weak cycling preference alongside the e-bike uptake and competing PT attractiveness associated with below-benchmark outcomes.

The baseline box contains 625 cases, all below-benchmark: density is 100\%, coverage 62.63\%, and mass 31.25\%. It is retained as a concentrated diagnostic core. Step~4 offers broader monitoring coverage (100\%) at 76.07\% density and 65.60\% mass. The selected East-stretch box also has 100\% density and 31.25\% mass, with 63.97\% coverage; the full-route box has 100\% density, 58.97\% coverage and 28.10\% mass. Complete narrowed ranges and the peeling trade-off are provided in Appendix~\ref{app:nb03-prim-evidence}. These in-sample regions do not guarantee performance in unseen futures. The saved surrogate's OD-delay accuracy warning and clipping at its training-domain boundaries limit interpretation.
